\begin{onehalfspace}

Questa sezione approfondisce la Sezione~\ref{sec:finale_5_3_8} del
Capitolo~\ref{chap:risultati}. Nel corpo restano il confronto finale fra
Gross code e surface code e il ruolo del decoder; qui sono conservati i
controlli di costruzione, gli sweep del decoder, la spiegazione del loro
comportamento, la validazione della ricostruzione a basso rumore e i
risultati sull'intera famiglia di codici. Tutti i numeri provengono dagli
artefatti della campagna M16, elencati in fondo alla sezione.

\subsection{Modello, metrica e criteri fissati in anticipo}
\label{app:gross_modello}

Il modello è code-capacity: in ogni shot ciascun qubit di dato subisce un
errore $X$ con probabilità $q$, indipendentemente dagli altri, e la
sindrome è misurata senza errori in un solo ciclo. Per la struttura CSS
dei codici considerati il settore $Z$ è equivalente e non viene simulato.
La metrica è la probabilità che almeno uno dei $k$ qubit logici del blocco
fallisca; per il surface code, che protegge un qubit logico per patch, si
usano $k$ patch indipendenti e $P_{\text{blocco}} = 1-(1-p_L)^k$.

Prima degli approfondimenti sono stati fissati tre criteri di
adeguatezza, pensati per rendere il confronto statisticamente solido e
non per favorire uno dei due codici:
\begin{enumerate}
\item ogni punto del confronto deve avere almeno 100 fallimenti; in caso
  contrario si riporta una stima a bassa statistica oppure, con zero
  fallimenti su $N$ shot, il limite superiore al 95\% $3/N$;
\item dove il primo criterio è soddisfatto, l'intervallo di confidenza al
  95\% deve avere larghezza relativa entro il 20\%;
\item le varianti ragionevoli del decoder devono dare risultati
  compatibili con il riferimento; se una variante lo migliora di oltre il
  10\% con p di McNemar inferiore a 0,01, diventa il nuovo riferimento.
\end{enumerate}
Un esito è dichiarato ``migliore'' o ``peggiore'' solo con intervalli al
95\% non sovrapposti. Il modello, la metrica e i punti del confronto non
sono stati modificati dopo aver visto i risultati.

\subsection{Controlli di costruzione e distanze}
\label{app:gross_controlli}

Per ciascun codice lo script verifica $n$, $k$ e i ranghi delle matrici
di parità; per il Gross code $n = 144$, $k = 12$, rango 66 per $H_X$ e
$H_Z$ e controlli di peso 6, come nella costruzione di
riferimento~\cite{bravyi2024}. La distanza è cercata con un metodo a
insiemi d'informazione casuali: si permutano le colonne, si porta la
matrice in forma ridotta e si leggono operatori logici di peso basso,
verificandoli poi in modo indipendente. Il metodo fornisce limiti
superiori alla distanza: trovare un logico del peso pubblicato conferma
la costruzione, trovarne uno più leggero la invaliderebbe.

\begin{table}[H]
\centering
\small
\begin{tabular}{@{}lccr@{}}
\toprule
\textbf{Codice} & \textbf{Distanza attesa} & \textbf{Trovata ($X$ / $Z$)} & \textbf{Tempo} \\
\midrule
surface $d = 5, 7, 9$ & 5, 7, 9 & 5, 7, 9 & 2--26 s \\
$[[72,12,6]]$ & 6 & 6 / 6 & 26 s \\
$[[90,8,10]]$ & 10 & 10 / 10 & 44 s \\
$[[108,8,10]]$ & 10 & 10 / 10 & 79 s \\
$[[144,12,12]]$ & 12 & 12 / 12 & 228 s \\
$[[288,12,18]]$ & 18 & 18 / 18 & 2427 s \\
\bottomrule
\end{tabular}
\caption[Distanze dei codici bivariate bicycle]{Verifica delle distanze
con 40\,000 iterazioni per settore. Per le distanze 3, 5 e 7 il surface
code è stato controllato anche per forza bruta.}
\label{tab:app_gross_distanze}
\end{table}

In nessuna corsa una correzione ha violato la sindrome, e sul surface
code \ac{MWPM} e \ac{BPOSD} danno lo stesso $p_L$ entro l'incertezza a
tutte le distanze (per esempio $4{,}55\times10^{-3}$ contro
$4{,}43\times10^{-3}$ a $d = 13$, $q = 5\%$): il riferimento non è
penalizzato dal proprio decoder.

\subsection{Scelta del decoder}
\label{app:gross_decoder}

Le prime corse usavano \ac{BPOSD} con belief propagation product-sum e
poi min-sum con fattore di scala 0,625, entrambe con aggiornamento
parallelo dei messaggi. Con queste configurazioni il Gross code
risultava peggiore di 12 patch a distanza 9 per $q \le 1\%$. Poiché il
terzo criterio non era soddisfatto, sono stati eseguiti due sweep sugli
stessi shot. La Tabella~\ref{tab:app_gross_sweep} riporta le varianti
più informative del primo.

\begin{table}[H]
\centering
\footnotesize
\begin{tabular}{@{}lcccc@{}}
\toprule
\textbf{Variante} & $q = 0{,}5\%$ & $1\%$ & $2\%$ & $3\%$ \\
\midrule
rif.\ min-sum 0,625 parallelo & $9{,}3\times10^{-6}$ & $9{,}2\times10^{-5}$ & $7{,}7\times10^{-4}$ & $2{,}95\times10^{-3}$ \\
seriale (0,625) & 0 & $5\times10^{-7}$ & $8{,}0\times10^{-5}$ & $1{,}15\times10^{-3}$ \\
min-sum 0,5, parallelo & 0 & $5\times10^{-7}$ & $8{,}0\times10^{-5}$ & $1{,}20\times10^{-3}$ \\
OSD-CS ordine 60 & 0 & $6\times10^{-6}$ & $1{,}3\times10^{-4}$ & $1{,}3\times10^{-3}$ \\
1000 iterazioni & $2{,}5\times10^{-5}$ & $1{,}6\times10^{-4}$ & $8{,}9\times10^{-4}$ & $3{,}1\times10^{-3}$ \\
product-sum & $3{,}7\times10^{-5}$ & $2{,}9\times10^{-4}$ & $1{,}7\times10^{-3}$ & $4{,}6\times10^{-3}$ \\
\bottomrule
\end{tabular}
\caption[Sweep del decoder sul Gross code]{Primo sweep: $p_L$ del
blocco su $4$, $2$, $1$ e $0{,}4$ milioni di shot, identici per tutte le
varianti. Le altre sei varianti (fattori 0,75, 0,875 e 1,0, OSD-CS di
ordine 30, OSD-E di ordine 7, BP+LSD) sono nell'artefatto dello sweep;
nessuna è migliore della schedulazione seriale.}
\label{tab:app_gross_sweep}
\end{table}

La schedulazione seriale e il fattore 0,5 riducono $p_L$ di un ordine di
grandezza o più a $q = 1$--$2\%$ e di circa 2,5 volte a $q = 3\%$, con p
di McNemar compresi fra $10^{-9}$ e $10^{-150}$. Il secondo sweep, con
riferimento seriale e undici varianti attorno ad esso su 4, 2, 0,8 e 0,4
milioni di shot a $q = 1$, 2, 3 e 4\%, non ha trovato miglioramenti
superiori al 10\%: le varianti migliori, con OSD di ordine 60, guadagnano
fra il 3 e il 5\%. A $q = 1\%$ nove decoder diversi sbagliano gli stessi
cinque shot. Il decoder seriale è quindi il riferimento definitivo; la
scelta, fatta sugli stessi dati del confronto, è stata confermata su semi
indipendenti, con cui il Gross code concorda entro il 10\%.

\paragraph{Perché la schedulazione conta.} Un terzo esperimento ha
variato il fattore di scala del min-sum parallelo da 0,400 a 1,000 a
passi di 0,025, su $5\times10^5$ shot identici a $q = 2\%$ e $3\%$,
registrando per ogni fallimento se la belief propagation aveva
convergito. Nelle configurazioni peggiori l'89--98\% dei fallimenti
avviene quando la belief propagation non converge e l'OSD sceglie la
correzione sbagliata; la convergenza a una soluzione errata resta
minoritaria. Il fallimento logico si scompone quindi in
$p_L \approx P(\text{non convergenza}) \times P(\text{errore OSD}\mid\text{non convergenza})$.
Al crescere del fattore il primo termine diminuisce, ma il secondo sale
bruscamente fra 0,5 e 0,525 e arriva al 10--20\%: nei casi difficili le
probabilità che la belief propagation passa all'OSD diventano
fuorvianti. Il prodotto ha un massimo intorno a 0,7, il che spiega
l'andamento apparentemente irregolare della griglia grossolana del primo
sweep. Anche fermare la belief propagation dopo 10 iterazioni invece che
dopo 100 riduce gli errori di circa tre volte, coerentemente con questa
interpretazione. Con la schedulazione seriale la belief propagation non
converge nello 0,07\% dei casi contro lo 0,77\% del parallelo, e $p_L$
resta piatto fra i fattori 0,4 e 0,9. Il parallelo nel suo punto
migliore e il seriale raggiungono lo stesso livello, segno che quel
livello appartiene al codice più che ai decoder.

\subsection{Basso rumore: decomposizione per peso e validazione}
\label{app:gross_peso}

A $q = 0{,}5\%$ il Gross code con decoder seriale fallisce circa una volta
ogni $10^8$ shot, e una stima diretta richiederebbe circa $10^{10}$ shot.
Poiché con errori indipendenti e priore uniforme le decisioni del decoder
non dipendono da $q$, vale esattamente
\begin{equation}
p_L(q) = \sum_{w=0}^{n} \binom{n}{w} q^{w} (1-q)^{n-w} f(w),
\label{eq:app_gross_peso}
\end{equation}
dove $f(w)$ è la probabilità di fallimento con esattamente $w$ errori.
Per il Gross code i pesi da 0 a 5 sono stati enumerati in modo
esaustivo: tutte le 498\,685\,189 configurazioni sono state decodificate
senza alcun fallimento, in circa 40 minuti. Per il surface code \ac{MWPM}
corregge per costruzione ogni errore di peso non superiore a $(d-1)/2$.
I pesi successivi sono stati campionati fino a $2\times10^7$ prove o 1000
fallimenti per peso.

La ricostruzione è stata accettata solo dopo due controlli diretti
indipendenti a $q = 1\%$, con soglia di accordo del 20\%: per il Gross
code $9{,}4\times10^{-7}$ (94 fallimenti su $10^8$ shot) contro
$1{,}01\times10^{-6}$ ricostruito, scarto del 7\%; per il surface code a
distanza 13 $1{,}350\times10^{-7}$ (135 su $10^9$) contro
$1{,}357\times10^{-7}$, scarto dello 0,5\%. I residui logici più leggeri
osservati hanno peso 12, 11 e 13 per il Gross code e per il surface code
a distanza 11 e 13, pari alle rispettive distanze.

\subsection{Famiglia bivariate bicycle e distanze maggiori}
\label{app:gross_famiglia}

La Tabella~\ref{tab:app_gross_famiglia} riporta i cinque codici della
famiglia, simulati con il decoder seriale e semi nuovi fino a $10^7$
shot o 300 fallimenti per punto.

\begin{table}[H]
\centering
\footnotesize
\setlength{\tabcolsep}{3pt}
\begin{tabular}{@{}lcccccl@{}}
\toprule
\textbf{Codice} & $0{,}5\%$ & $1\%$ & $2\%$ & $3\%$ & $5\%$ & \textbf{$d$ equiv.} \\
\midrule
$[[72,12,6]]$ & $1{,}1\times10^{-4}$ & $1{,}0\times10^{-3}$ & $9{,}2\times10^{-3}$ & $3{,}3\times10^{-2}$ & $0{,}15$ & 5--7 \\
$[[90,8,10]]$ & $10^{-7}$ (1 ev.) & $1{,}0\times10^{-5}$ & $3{,}7\times10^{-4}$ & $3{,}3\times10^{-3}$ & $4{,}2\times10^{-2}$ & 9--13 \\
$[[108,8,10]]$ & 0 su $10^7$ & $1{,}9\times10^{-6}$ & $1{,}5\times10^{-4}$ & $1{,}7\times10^{-3}$ & $2{,}9\times10^{-2}$ & $>11$ \\
$[[144,12,12]]$ & 0 su $10^7$ & $1{,}1\times10^{-6}$ & $8{,}2\times10^{-5}$ & $1{,}0\times10^{-3}$ & $2{,}8\times10^{-2}$ & $>13$ da 2\% \\
$[[288,12,18]]$ & 0 su $10^7$ & 0 su $10^7$ & $10^{-7}$ (1 ev.) & $9\times10^{-6}$ & $3{,}7\times10^{-3}$ & $>13$ da 1\% \\
\bottomrule
\end{tabular}
\caption[Famiglia bivariate bicycle]{Probabilità di fallimento del
blocco da $k$ qubit logici per i cinque codici bivariate bicycle, con il
decoder seriale. L'ultima colonna indica la distanza del surface code
le cui $k$ patch hanno lo stesso fallimento del blocco.}
\label{tab:app_gross_famiglia}
\end{table}

Tutti i codici con distanza almeno 10 battono le patch di surface code
che usano lo stesso numero di qubit di dato o più. A parità di $k$ e $d$,
$[[108,8,10]]$ fa circa cinque volte meglio di $[[90,8,10]]$ a
$q = 1\%$: la distanza non esaurisce la descrizione del codice. Per
estendere il confronto del codice più grande, il surface code è stato
simulato anche a distanza 15 e 17, fino a $10^8$ shot per punto, e
$[[288,12,18]]$ fino a $10^8$ shot a $q = 2\%$ e $3\%$. Da $q = 2\%$ in su
$[[288,12,18]]$ batte 12 patch fino a distanza 17, che usano 3468 qubit
di dato contro 288: $5{,}0\times10^{-8}$ contro $1{,}28\times10^{-5}$ a
$q = 2\%$ e $9{,}5\times10^{-6}$ contro $4{,}2\times10^{-4}$ a $q = 3\%$.
A $q = 1\%$ batte le patch fino a distanza 13; oltre, e a $q = 0{,}5\%$,
i fallimenti osservati non bastano a decidere.

\subsection{Artefatti e limiti}
\label{app:gross_artefatti}

Gli script e gli artefatti sono in
\texttt{Extra/experiments/M16\_qldpc\_gross}; ogni file JSON segue lo
schema 2.0 con manifest, seed, SHA-256 degli script e versioni software
(NumPy 2.5.1, SciPy 1.18.0, ldpc 2.4.1, PyMatching 2.4.0, Stim 1.16.0,
Python 3.12.3).

\begin{table}[H]
\centering
\footnotesize
\begin{tabularx}{\textwidth}{lX}
\toprule
\textbf{Risultato} & \textbf{Artefatto in \texttt{artifacts/v2\_\allowbreak{}20260926}} \\
\midrule
sweep del decoder & \texttt{results\_\allowbreak{}M16\_\allowbreak{}test1\_\allowbreak{}sweep\_\allowbreak{}decoder\_\allowbreak{}20260926\_\allowbreak{}104516.json}, \texttt{\dots\_\allowbreak{}giro1b\_\allowbreak{}20260926\_\allowbreak{}105429.json} \\
$f(w)$ & \texttt{results\_\allowbreak{}M16\_\allowbreak{}test2\_\allowbreak{}fallimenti\_\allowbreak{}20260926\_\allowbreak{}121335.json} \\
distanze & \texttt{results\_\allowbreak{}M16\_\allowbreak{}test3\_\allowbreak{}distanza\_\allowbreak{}20260926\_\allowbreak{}115305.json} \\
famiglia BB & \texttt{results\_\allowbreak{}M16\_\allowbreak{}test4\_\allowbreak{}famiglia\_\allowbreak{}bb\_\allowbreak{}20260926\_\allowbreak{}120439.json} \\
meccanismo del decoder & \texttt{results\_\allowbreak{}M16\_\allowbreak{}anomalia\_\allowbreak{}scala\_\allowbreak{}20260926\_\allowbreak{}122400.json} \\
enumerazione esaustiva & \texttt{results\_\allowbreak{}M16\_\allowbreak{}esaustivo\_\allowbreak{}gross\_\allowbreak{}144\_\allowbreak{}12\_\allowbreak{}12\_\allowbreak{}serial\_\allowbreak{}w5\_\allowbreak{}20260926\_\allowbreak{}133538.json} \\
ricostruzione a basso rumore & \texttt{results\_\allowbreak{}M16\_\allowbreak{}test2\_\allowbreak{}fallimenti\_\allowbreak{}basso\_\allowbreak{}rumore\_\allowbreak{}20260926\_\allowbreak{}134543.json} \\
controlli diretti & \texttt{results\_\allowbreak{}M16\_\allowbreak{}test4\_\allowbreak{}famiglia\_\allowbreak{}bb\_\allowbreak{}gross\_\allowbreak{}q001\_\allowbreak{}1e8\_\allowbreak{}20260926\_\allowbreak{}135157.json}, \texttt{results\_\allowbreak{}M16\_\allowbreak{}gross\_\allowbreak{}code\_\allowbreak{}capacity\_\allowbreak{}d13\_\allowbreak{}q001\_\allowbreak{}20260926\_\allowbreak{}143129.json} \\
confronto a basso rumore & \texttt{analysis\_\allowbreak{}M16\_\allowbreak{}confronto\_\allowbreak{}basso\_\allowbreak{}rumore\_\allowbreak{}20260926\_\allowbreak{}143130.json} \\
confronto con decoder definitivo & \texttt{analysis\_\allowbreak{}M16\_\allowbreak{}consolidata\_\allowbreak{}serial\_\allowbreak{}20260926\_\allowbreak{}120628.json} \\
distanze 15 e 17 & \texttt{results\_\allowbreak{}M16\_\allowbreak{}gross\_\allowbreak{}code\_\allowbreak{}capacity\_\allowbreak{}surface\_\allowbreak{}d15\_\allowbreak{}d17\_\allowbreak{}20260926\_\allowbreak{}153636.json}, \texttt{analysis\_\allowbreak{}M16\_\allowbreak{}consolidata\_\allowbreak{}bb288\_\allowbreak{}d17\_\allowbreak{}20260926\_\allowbreak{}160258.json} \\
\bottomrule
\end{tabularx}
\caption[Artefatti della campagna M16]{Artefatti canonici della campagna
M16. Le corse con decoder parallelo restano nella stessa cartella come
traccia e non sono usate nel confronto finale.}
\label{tab:app_gross_artefatti}
\end{table}

Restano i limiti del modello: niente errori di misura né circuito di
estrazione delle sindromi, errori $X$ indipendenti, distanze verificate
come limiti superiori e nessun confronto a livello di circuito con il
surface code della Sezione~\ref{sec:finale_5_2_3} o con i risultati
pubblicati da IBM. Il confronto conta inoltre i soli qubit di dato:
includendo le ancille, il Gross code usa 288 qubit e 12 patch a distanza
13 ne usano 4044.

\end{onehalfspace}
